\section{Эмпирический анализ сопутствующих индикаторов и показателей промышленного цикла}

\subsection{Политэкономические показатели}
Динамика органического строения капитала и нормы прибыли отражает внутренне взаимосвязанное движение капиталистического накопления в рамках промышленного цикла: развитие производительных сил изменяет соотношение между овеществленным и живым трудом, вследствие чего рост капиталовооруженности производства сочетается с изменением условий создания и реализации прибавочной стоимости.

Органическое строение капитала имеет долгосрочную тенденцию к повышению вследствие механизации, автоматизации и научно-технического прогресса, увеличивающих массу средств производства, приходящуюся на одного работника; в фазе подъема, особенно перед кризисом, его рост ускоряется из-за ввода производственных мощностей, созданных в результате предшествующих массовых инвестиций, а во время кризиса дополнительно усиливается вследствие резкого сокращения занятости и фонда заработной платы при сравнительно медленном уменьшении уже накопленного основного капитала.

Норма прибыли, как правило, движется в направлении, обратном органическому строению капитала, и обнаруживает долгосрочную тенденцию к понижению: примерно с середины подъема она начинает снижаться по мере того, как накопление постоянного капитала опережает рост прибавочной стоимости, а в период кризиса падение усиливается вследствие сокращения производства, недозагрузки мощностей, банкротств и обесценения части капитала. При этом положительная агрегированная норма прибыли означает лишь положительный итог для совокупного общественного капитала и не предполагает прибыльности каждого предприятия или новых инвестиций, поскольку убытки множества мелких фирм могут перекрываться высокой совокупной прибылью нескольких крупных корпораций.

Таким образом, повышение органического строения капитала в ходе подъема создает предпосылки снижения прибыльности и последующего кризиса, тогда как сам кризис посредством обесценения и разрушения части капитала, сокращения избыточных мощностей и удешевления рабочей силы временно восстанавливает условия прибыльного накопления, одновременно подготавливая начало нового промышленного цикла. $\dagger$ (Приложение \ref{app:emp}, стр. \pageref{fig:occ-profit})

\subsection{Промышленное производство}
Динамика показателей промышленного производства в рамках промышленного цикла отражает последовательное изменение условий накопления, загрузки и обновления производственного аппарата, вследствие чего колебания инвестиций, заказов, возраста основного капитала, использования мощностей и запасов образуют взаимосвязанную систему, позволяющую проследить переход экономики от оживления и подъема к кризису и последующему восстановлению.

Средний возраст основного капитала начинает снижаться в начале фазы подъема, когда расширение спроса и улучшение условий прибыльного накопления стимулируют массовое обновление оборудования, ввод новых производственных мощностей и замену физически либо морально устаревших средств труда; в период кризиса и последующего оживления данный показатель, напротив, возрастает, поскольку предприятия ограничивают крупные капиталовложения, а восстановление выпуска осуществляется преимущественно за счет повторной загрузки уже существующего производственного аппарата.

Инвестиции в основной капитал начинают расти с переходом к фазе подъема, когда увеличение загрузки мощностей, прибыли и ожидаемого спроса побуждает предприятия расширять и модернизировать производство, однако ближе к завершению подъема их динамика замедляется и сменяется снижением вследствие ухудшения прибыльности, накопления избыточных мощностей и ослабления инвестиционных ожиданий; на фазе оживления инвестиции, как правило, стагнируют или увеличиваются незначительно, поскольку первоначальное восстановление производства происходит без масштабного расширения основного капитала.

Портфель заказов в целом повторяет динамику инвестиций в основной капитал, поскольку расширение производственных планов и инвестиционной активности сопровождается ростом спроса на оборудование, материалы и промежуточную продукцию, тогда как замедление накопления в конце подъема и сокращение производства в период кризиса приводят к уменьшению числа новых заказов, отмене ранее заключенных контрактов и ухудшению загрузки предприятий, выпускающих средства производства.

Загрузка производственных мощностей изменяется практически одновременно с объемом промышленного производства: она повышается в ходе оживления и подъема по мере вовлечения простаивающего оборудования и расширения выпуска, достигает максимальных значений вблизи вершины цикла, а затем резко снижается во время кризиса, когда сокращение спроса и производства превращает часть накопленного основного капитала в бездействующие или недоиспользуемые мощности.

Товарно-материальные запасы сокращаются в период кризиса вследствие снижения производства, распродажи ранее накопленной продукции и стремления предприятий высвободить связанный в запасах оборотный капитал; в ходе последующего подъема темпы их накопления ускоряются примерно к середине фазы под воздействием расширения выпуска и ожиданий дальнейшего роста спроса, однако ближе к ее завершению замедляются, поскольку ухудшение реализации, рост издержек хранения и признаки приближающегося перепроизводства заставляют предприятия осторожнее формировать новые запасы.

Таким образом, совместная динамика среднего возраста основного капитала, инвестиций, портфеля заказов, загрузки производственных мощностей и товарно-материальных запасов позволяет рассматривать промышленный цикл как целостный процесс последовательного расширения, перенакопления, кризисного сокращения и последующего восстановления производства; при этом опережающее ухудшение инвестиционной активности и заказов, замедление накопления запасов и последующее снижение загрузки мощностей выступают важными признаками перехода от подъема к кризису, тогда как оживление начинается преимущественно с более полного использования уже существующего производственного аппарата и лишь затем перерастает в новое расширение основного капитала. $\dagger$ (Приложение \ref{app:emp}, стр. \pageref{fig:average-age}-\pageref{fig:inventories})

\subsection{Занятость и уровень цен}
Динамика цен и безработицы в рамках промышленного цикла отражает изменения соотношения между совокупным спросом, масштабами производства, загрузкой производственных мощностей и состоянием рынка труда, вследствие чего инфляционные и дефляционные процессы, а также колебания занятости выступают важными характеристиками перехода экономики от кризиса к оживлению и подъёму.

В начале фазы кризиса темпы прироста индекса цен производителей и индекса потребительских цен могут временно увеличиваться под воздействием накопившегося ранее роста издержек, инерции контрактных цен и запаздывающей реакции рынка на ухудшение конъюнктуры, однако по мере углубления спада, сокращения спроса и реализации избыточных товарных запасов они резко снижаются, а в отдельных случаях становятся отрицательными; на фазе оживления темпы прироста цен постепенно восстанавливаются и в ходе подъёма достигают максимальных значений, причём индекс цен производителей обычно отличается большей волатильностью, поскольку быстрее реагирует на изменения стоимости сырья, энергии и промежуточной продукции, тогда как индекс потребительских цен изменяется более устойчиво и значительно реже демонстрирует отрицательную динамику.

Безработица начинает быстро расти в период кризиса вследствие сокращения производства, остановки предприятий, банкротств и массового высвобождения рабочей силы, однако на фазе оживления темпы её увеличения замедляются по мере восстановления выпуска и повторного вовлечения работников в производство; с переходом к фазе подъёма уровень безработицы начинает устойчиво снижаться и к её завершению достигает минимальных значений, отражая высокую степень загрузки рабочей силы и относительное сокращение резервной армии труда.

Таким образом, ценовая динамика и безработица движутся в промышленном цикле не изолированно, а как взаимосвязанные проявления изменения условий воспроизводства: кризис сопровождается ослаблением ценового давления и ростом безработицы, оживление — восстановлением темпов роста цен и стабилизацией рынка труда, а подъём — ускорением инфляции и сокращением безработицы до минимального уровня, что одновременно выражает расширение производства и постепенное исчерпание свободных производственных и трудовых ресурсов. $\dagger$ (Приложение \ref{app:emp}, стр. \pageref{fig:prices}-\pageref{fig:unemployment})

\subsection{Денежное обращение и фондовый рынок}
Динамика денежного обращения и фондовых индексов в рамках промышленного цикла отражает изменения масштабов производства, кредитной активности, ликвидности и ожиданий экономических субъектов, вследствие чего колебания денежной массы, скорости обращения денег и рыночной стоимости финансовых активов образуют взаимосвязанную систему, реагирующую как на текущее состояние воспроизводства, так и на предполагаемое направление его дальнейшего движения.

Во время фазы кризиса реальный объем денежной массы может сокращаться вследствие сжатия банковского кредитования, роста неплатежей, погашения долгов и общего снижения деловой активности либо, напротив, увеличиваться, если государство и центральный банк проводят стимулирующую антикризисную политику посредством расширения денежного предложения и поддержки финансовой системы; на фазах оживления и подъема денежная масса, как правило, возрастает под воздействием восстановления кредитования, роста доходов и расширения хозяйственного оборота, хотя к концу подъема возможно ее некоторое сокращение или замедление темпов роста в случае ухудшения прибыльности, повышения кредитных рисков и уменьшения объемов новых займов.

Скорость обращения денег в целом характеризуется динамикой, обратной изменению объема денежной массы, поскольку расширение денежных остатков нередко сопровождается снижением интенсивности их использования, тогда как ограничение ликвидности может побуждать экономических субъектов ускорять расчеты; во время кризиса скорость обращения обычно падает вследствие сокращения производства, инвестиций и потребительских расходов, накопления ликвидных резервов и нарушения платежных цепочек, однако на фазах оживления и подъема ее движение может различаться в зависимости от соотношения темпов роста денежной массы, номинального выпуска, кредитной активности и предпочтения ликвидности.

Фондовые индексы в целом повторяют направление динамики промышленного производства, повышаясь в периоды расширения и снижаясь при ухудшении экономической конъюнктуры, однако обычно обладают опережающим характером, поскольку биржевые цены формируются на основе ожиданий будущих прибылей, процентных ставок и масштабов деловой активности; вследствие этого их падение нередко начинается еще до явного сокращения промышленного выпуска, тогда как восстановление котировок может происходить уже в конце кризиса или на ранней стадии оживления, когда фактические показатели производства еще остаются низкими.

Таким образом, изменения денежной массы, скорости обращения денег и фондовых индексов выражают различные стороны циклического движения капитала: денежное обращение отражает состояние кредита, ликвидности и хозяйственного оборота, тогда как фондовый рынок быстрее реагирует на ожидаемые изменения прибыльности и условий накопления, благодаря чему совместный анализ этих показателей позволяет точнее выявлять переходы между кризисом, оживлением и подъемом. $\dagger$ (Приложение \ref{app:emp}, стр. \pageref{fig:m2-velocity}-\pageref{fig:stock-market})

\subsection{Международная торговля и движение капитала}
Динамика чистого экспорта и прямых иностранных инвестиций в рамках промышленного цикла отражает изменение внутреннего спроса, инвестиционной активности и международного движения капитала, вследствие чего внешнеторговые потоки и трансграничные капиталовложения выступают важными показателями положения национальной экономики в мировой системе воспроизводства.

Для экономики США характерна долгосрочная тенденция к сохранению и расширению торгового дефицита, обусловленная устойчивым превышением импорта товаров и услуг над экспортом, однако в период кризиса чистый экспорт, как правило, увеличивается вследствие более резкого сокращения импорта под воздействием падения внутреннего потребления, инвестиций и спроса на промежуточную продукцию; при этом улучшение торгового баланса в данном случае выражает не укрепление экспортного потенциала, а прежде всего кризисное сжатие внутреннего рынка и хозяйственного оборота.

Динамика исходящих и входящих прямых иностранных инвестиций в целом повторяет движение частных инвестиций в основной капитал: их рост замедляется и сменяется спадом к завершению фазы подъема, когда ухудшаются ожидания прибыльности и возрастают риски новых капиталовложений, а во время кризиса сокращение усиливается вследствие падения спроса, дефицита кредитных ресурсов и общей неопределенности; на фазе оживления трансграничные инвестиционные потоки временно восстанавливаются по мере стабилизации производства и финансовой системы, после чего с началом нового подъема переходят к более устойчивому росту, сопровождающему расширение накопления капитала как внутри страны, так и за ее пределами.

Таким образом, чистый экспорт и прямые иностранные инвестиции реагируют на промышленный цикл различным образом: кризисное сокращение импорта способно временно улучшать торговый баланс, тогда как международные инвестиционные потоки обычно уменьшаются вместе с общей инвестиционной активностью, а их последующее восстановление свидетельствует о возобновлении накопления, укреплении деловых ожиданий и расширении международного движения капитала. $\dagger$ (Приложение \ref{app:emp}, стр. \pageref{fig:net-export}-\pageref{fig:fdi-flow})

\subsection{Государственная регулирующая политика}
Динамика государственных расходов и ключевой процентной ставки в рамках промышленного цикла отражает изменение направленности макроэкономической политики, посредством которой государство стремится воздействовать на совокупный спрос, инвестиционную активность, занятость и устойчивость финансовой системы на различных фазах экономической конъюнктуры.

Государственные расходы в экономике США имеют долгосрочную тенденцию к увеличению под воздействием расширения масштабов государственного регулирования, социальных обязательств, военных затрат и финансирования инфраструктуры, однако в период кризиса темпы их прироста обычно ускоряются вследствие проведения стимулирующей налогово-бюджетной политики, включающей рост государственных закупок, трансфертов, субсидий и антикризисных программ, направленных на частичную компенсацию сокращения частного потребления и инвестиций.

Ключевая процентная ставка, как правило, повышается ближе к завершению фазы подъема, когда центральный банк стремится ограничить инфляционное давление, чрезмерное кредитное расширение и перегрев экономики, тогда как во время кризиса она снижается с целью удешевления заемных средств, поддержки банковской ликвидности, стимулирования кредитования и предотвращения дальнейшего сокращения производства и занятости; таким образом, ее циклическая динамика соответствует переходу от преимущественно сдерживающей денежно-кредитной политики в конце подъема к стимулирующей политике в период спада.

Таким образом, государственные расходы и ключевая ставка выступают важнейшими инструментами антициклического регулирования: в период кризиса расширение бюджетных расходов и смягчение денежно-кредитных условий используются для поддержания совокупного спроса и стабилизации воспроизводства, тогда как к концу подъема повышение ставки и более умеренная бюджетная динамика направлены на ограничение инфляции и накопления финансовых диспропорций.

Вместе с тем государственная налогово-бюджетная и денежно-кредитная политика способна лишь смягчать или, напротив, усиливать отдельные проявления кризиса, изменять его продолжительность, глубину и конкретную форму, но не устранять циклический характер капиталистического воспроизводства. Это обусловлено тем, что государственное регулирование не преодолевает имманентные противоречия капиталистического способа производства – между общественным характером производства и частной формой присвоения, производством и ограниченным платежеспособным спросом, накоплением капитала и тенденцией нормы прибыли к понижению, – которые образуют фундаментальную основу периодического возникновения кризисов и промышленного цикла. $\dagger$ (Приложение \ref{app:emp}, стр. \pageref{fig:gov-consumption-key-rate})